This appendix documents the protocol used to query the raw measurement
database from which the search, footprint, stability, and construction results
reported in this thesis are derived. The end-to-end wall-clock times of
Table~\ref{tab:e2e} were collected in a dedicated pipeline benchmark conducted
on 2026-05-22 and are documented in the companion repository's synthesis report
(\texttt{parallel-aho-corasick/docs/tcc-synthesis.html}, Section~3.9); the raw
sweep script and timing log are not currently archived in the repository.
% REVIEW-TODO: archive scripts/run_f3_v3flat_build_par_sweep.sh and f3_sweep.log for full reproducibility
This allows the primary figures to be reproduced and audited independently.
The benchmark campaign stores its measurements in a single typed SQLite
database located at \texttt{parallel-aho-corasick/runs/overnight/sweep.db},
containing one \texttt{runs} table (one row per timed execution) and a
set of aggregating views. The headline figures are retrieved from these views
instead of the raw execution logs.

\subsecao{Database Schema}

Each row of the \texttt{runs} table identifies an execution by its phase,
pattern dictionary, corpus, searcher strategy, and thread count. It records the
timing statistics (minimum, mean, median, and maximum wall-clock time in
milliseconds, along with the coefficient of variation across the timed
iterations), the derived throughput (in MB/s and Gbit/s), the input volume and
match count, and the automaton descriptors (pattern count, number of states,
footprint in kibibytes, and sequential or parallel build time). The views
listed in Table~\ref{tab:dbviews} pre-aggregate the data so no headline
figure requires ad hoc post-processing of the raw rows.

\begin{table}[htbp]
    \centering
    \caption{Aggregating views of the measurement database.}
    \label{tab:dbviews}
    \small
    \begin{tabular}{@{} l p{8.8cm} @{}}
        \toprule
        \textbf{View} & \textbf{Question it answers} \\
        \midrule
        \texttt{v\_speedup}       & Speedup of each searcher against the single-threaded sequential baseline. \\ \addlinespace
        \texttt{v\_self\_speedup} & Self-relative scalability: throughput at $T$ threads over throughput at one thread. \\ \addlinespace
        \texttt{v\_footprint}     & Throughput as a function of the automaton footprint, at one and twelve threads. \\ \addlinespace
        \texttt{v\_build}         & Speedup of the parallel automaton construction over the sequential build. \\ \addlinespace
        \texttt{v\_best}          & Best-performing searcher for each dictionary, corpus and thread count. \\ \addlinespace
        \texttt{v\_correctness}   & Number of distinct match counts per dictionary and corpus (a value of one indicates agreement). \\
        \bottomrule
    \end{tabular}

    {\footnotesize\raggedright \textbf{Source:} Prepared by the author. \par}
\end{table}

\subsecao{Canonical Queries}

The queries below reproduce the main results of the study. They can be executed
against the database using the standard \texttt{sqlite3} client.

\begin{verbatim}
-- Speedup curve of a searcher (Snort dictionary, canonical corpus)
SELECT thr, speedup_vs_seq FROM v_speedup
WHERE patterns='patterns_snort' AND corpus='enron_corpus'
  AND searcher='pthread_dynamic' ORDER BY thr;

-- Peak speedup per scenario (winning searcher and thread count)
SELECT patterns, corpus, searcher, thr, speedup_vs_seq FROM v_speedup w
WHERE speedup_vs_seq = (SELECT MAX(speedup_vs_seq) FROM v_speedup x
  WHERE x.patterns = w.patterns AND x.corpus = w.corpus)
ORDER BY patterns, corpus;

-- Footprint versus throughput at twelve threads (cache blowout)
SELECT patterns, automaton_mib, searcher, mbps FROM v_footprint
WHERE thr = 12 ORDER BY automaton_mib;

-- Speedup of the parallel automaton construction
SELECT patterns, build_threads, build_speedup FROM v_build
ORDER BY patterns, build_threads;
\end{verbatim}

Before any timing was reported, the correctness of every parallel strategy was
verified using the \texttt{v\_correctness} view. This view returns a single
distinct match count for each dictionary and corpus, confirming that all
strategies reproduce the exact match set of the sequential baseline.
